\everymath{\displaystyle}

%%%%%%%%%%%%%%%%%%%%%%%%%%%%%%%%%%%%%%%%
%           main body
%%%%%%%%%%%%%%%%%%%%%%%%%%%%%%%%%%%%%%%%

%-----------------------------------
%	TITLE PAGE
%-----------------------------------

\title[第二章\\
勒贝格测度]{\texorpdfstring{\LARGE \bfseries 第二章\quad 勒贝格测度}{第二章 勒贝格测度}}
\author[\S 2.2\\
外内测度-可测集]
{\Large \bfseries \S 2.2 有界点集的外内测度-可测集}
\date{}

%------------------------------------------------

\begin{document}

%------------------------------------------------

\begin{frame}

\titlepage

\end{frame}

%------------------------------------------------

\section{引言}

%------------------------------------------------

\begin{frame}{从 Riemann 积分到 Lebesgue 测度}

{
\larger[1]
Riemann 积分 $\int_a^b f(x)\,\mathrm{d}x = \lim\limits_{\|T\| \to 0} \sum\limits_{i=1}^n f(\xi_i) \Delta x_i$ 的核心思想是对\textbf{定义域}进行分割. 这一做法将简单的工作留给了积分区间的划分 $\Delta x_i$, 而函数值 $f(\xi_i)$ 本身可能相当复杂.
}

\pause

这种做法带来两个根本性的局限:

\pause

\begin{itemize}
\item \textbf{可积函数类狭窄}: 由 Lebesgue 判别法, Riemann 可积函数必定几乎处处连续. 换言之, Riemann 积分只能处理``近似连续''的函数.
\pause
\item \textbf{极限运算不封闭}: Riemann 可积函数列的一致极限未必 Riemann 可积.
\end{itemize}

\pause

\begin{eg}[Dirichlet 函数]
$D(x) = \begin{cases} 1, & x \in \mathbb{Q}, \\ 0, & x \notin \mathbb{Q}. \end{cases}$ 在 $[0,1]$ 上处处不连续, 故不 Riemann 可积. 但从``测度''的角度看, 有理数集是``可以忽略''的, 其 Lebesgue 积分应为 $0$.
\end{eg}

\end{frame}

%------------------------------------------------

\begin{frame}{测度应满足的基本公理}

Lebesgue 的思想是反过来: 将函数值 $f(\xi_i)$ 简单化 (按值域分割), 而把复杂的工作留给定义域``长度''的推广. 这要求我们对 $\mathbb{R}$ 中的集合 $E$ 赋予一个``长度'' $m E$.

\pause

\begin{Def}[理想测度的公理]
设 $\mathscr{M}$ 是 $\mathbb{R}$ 的某些子集构成的集类, $m: \mathscr{M} \to [0, +\infty]$. 一个理想的测度应满足:
\begin{enumerate}
\item \textbf{非负性}: $m E \geqslant 0,$ $m \emptyset = 0;$
\item \textbf{规范性}: $m([0,1]) = 1;$
\item \textbf{平移不变性}: $E \in \mathscr{M} \Rightarrow E + x \in \mathscr{M},$ 且 $m(E + x) = m E$ ($\forall\, x \in \mathbb{R}$);
\item \textbf{完全可加性}: 若 $E_1, E_2, \ldots \in \mathscr{M}$ 互不相交, 则 $\bigcup\limits_{n=1}^\infty E_n \in \mathscr{M},$ 且 $m\left(\bigcup\limits_{n=1}^\infty E_n\right) = \sum\limits_{n=1}^\infty m E_n.$
\end{enumerate}
\end{Def}

\pause

\begin{question}
能否在 $\mathbb{R}$ 的\textbf{所有}子集上定义满足以上 4 条的测度?
\end{question}

\pause

\begin{attnblock}
答案是否定的 (参见补充材料中关于不可测集的讨论). 因此必须挑选出``好的''子集构成可测集类 $\mathscr{M}$. 本节的任务就是给出 $\mathscr{M}$ 的刻画并研究其初步性质.
\end{attnblock}

\end{frame}

%------------------------------------------------

\section{有界集的外测度}

%------------------------------------------------

\begin{frame}{外测度的定义}

\begin{Def}[外测度 (outer measure)]
设 $E \subset \mathbb{R}$ 为有界集. 称
\[
m^* E = \inf\left\{\sum_{n=1}^\infty l(I_n) : E \subset \bigcup_{n=1}^\infty I_n,\ I_n \text{ 为开区间}\right\}
\]
为 $E$ 的\alert{外测度 (outer measure)}. 其中 $l((a,b)) = b - a$ 表示区间长度.
\end{Def}

\pause

\begin{itemize}
\item 基本思想: 用可列个开区间从外部``覆盖'' $E$, 取下确界作为``长度''的逼近.
\item 由于 $E$ 有界, 存在有限区间包含它, 故下确界存在且 $0 \leqslant m^* E < +\infty$.
\end{itemize}

\pause

\begin{remark}
\begin{itemize}
\item 若将定义中的开区间改为闭区间或半开半闭区间, 得到的外测度相同.
\item 外测度是区间长度的推广: 对区间 $I$, 有 $m^* I = l(I)$ (需证明).
\end{itemize}
\end{remark}

\end{frame}

%------------------------------------------------

\begin{frame}{外测度的计算实例}

\begin{eg}[区间]
对 $E = [a,b], (a,b), [a,b), (a,b]$, 均有 $m^* E = b - a$ (证明见后).
\end{eg}

\pause

\begin{eg}[可列集]
$E = \mathbb{Q} \cap [0,1] = \{r_1, r_2, \ldots\}.$ 对任意 $\varepsilon > 0$, 取 $I_n = (r_n - \varepsilon/2^{n+1}, r_n + \varepsilon/2^{n+1})$, 则 $E \subset \bigcup_n I_n$, 且
\[
m^* E \leqslant \sum_{n=1}^\infty l(I_n) = \sum_{n=1}^\infty \frac{\varepsilon}{2^n} = \varepsilon.
\]
由 $\varepsilon$ 的任意性, $m^* E = 0$.
\end{eg}

\pause

\begin{eg}[Cantor 集]
$P_0 \subset [0,1]$ 为 Cantor 三分集. 由构造, $P_0 \subset F_n$, 而 $F_n$ 由 $2^n$ 个长度为 $3^{-n}$ 的闭区间组成, 故
\[
m^* P_0 \leqslant 2^n \cdot \frac{1}{3^n} = \left(\frac{2}{3}\right)^n \to 0 \quad (n \to \infty).
\]
从而 $m^* P_0 = 0$.
\end{eg}

\end{frame}

%------------------------------------------------

\section{外测度的基本性质}

%------------------------------------------------

\begin{frame}{外测度的基本性质 (1)}

\begin{prop}[基本性质]
设 $E, F$ 为有界集. 外测度 $m^*$ 满足:
\begin{enumerate}
\item \textbf{非负性}: $m^* E \geqslant 0,$ $m^* \emptyset = 0;$
\item \textbf{单调性}: 若 $E \subset F$, 则 $m^* E \leqslant m^* F;$
\item \textbf{平移不变性}: 对任意 $x \in \mathbb{R}$, 有 $m^*(E + x) = m^* E.$
\end{enumerate}
\end{prop}

\pause

\startproof[bg=yellow, fg=red]
\begin{enumerate}
\item 由定义, 外测度是非负数的下确界, 显然非负; 空集可被长度任意小的区间覆盖.
\pause
\item 覆盖 $F$ 的任一开区间列必覆盖 $E$, 故 $\{$覆盖 $F$ 的长度和$\} \subset \{$覆盖 $E$ 的长度和$\}$. 取下确界即得.
\pause
\item 将覆盖 $E$ 的区间列 $\{I_n\}$ 平移 $x$ 得到覆盖 $E+x$ 的区间列 $\{I_n + x\}$, 各区间长度不变, 故下确界相同.
\end{enumerate}
\qed

\end{frame}

%------------------------------------------------

\begin{frame}{外测度的基本性质 (2) — 次可加性}

\begin{prop}[次可加性 (subadditivity)]
对任意有界集列 $\{E_n\}_{n=1}^\infty$, 有
\[
m^*\left(\bigcup_{n=1}^\infty E_n\right) \leqslant \sum_{n=1}^\infty m^* E_n.
\]
\end{prop}

\pause

\startproof[bg=yellow, fg=red]
若右端为 $+\infty$, 不等式自然成立. 下设 $\sum_n m^* E_n < +\infty$.

\pause

对任意 $\varepsilon > 0$ 及每个 $n \in \mathbb{N}$, 由外测度定义, 存在开区间列 $\{I_{n,k}\}_{k=1}^\infty$ 使得
\[
E_n \subset \bigcup_{k=1}^\infty I_{n,k}, \qquad \sum_{k=1}^\infty l(I_{n,k}) < m^* E_n + \frac{\varepsilon}{2^n}.
\]

\pause

则 $\{I_{n,k}\}_{n,k \geqslant 1}$ 是可列个开区间, 覆盖 $\bigcup_{n=1}^\infty E_n$, 故
\[
m^*\left(\bigcup_{n=1}^\infty E_n\right) \leqslant \sum_{n=1}^\infty \sum_{k=1}^\infty l(I_{n,k}) < \sum_{n=1}^\infty \left(m^* E_n + \frac{\varepsilon}{2^n}\right) = \sum_{n=1}^\infty m^* E_n + \varepsilon.
\]
由 $\varepsilon$ 的任意性即得.
\end{frame}

%------------------------------------------------

\begin{frame}{外测度的基本性质 (3) — 次可加性}

\begin{attnblock}[等号一般不成立]
次可加性中的等号\textbf{一般不成立}. 这正是需要引入``可测集''概念的原因: 在可测集类上, 次可加性可以升级为可加性.
\end{attnblock}

\pause

\begin{eg}[次可加性取等号的例子]
若 $E_1, E_2, \ldots$ 互不相交且均为区间 (或开集), 则等号成立 (详见后续章节).
\end{eg}

\pause

\begin{remark}
次可加性是外测度最重要的性质之一. 它保证了外测度在某种意义下``不太大'', 是测度合理性的基本验证.
\end{remark}

\end{frame}

%------------------------------------------------

\begin{frame}{区间的外测度等于其长度}

\begin{thm}
对任意区间 $I$ (开、闭、半开半闭), 有 $m^* I = l(I).$
\end{thm}

\pause

\startproof[bg=yellow, fg=red]
设 $I = [a,b]$, 则 $l(I) = b - a$.

\pause

\textbf{先证 $m^* I \leqslant b - a$}: 对任意 $\varepsilon > 0$, 取 $I_1 = (a - \varepsilon, b + \varepsilon)$, $I_n = \emptyset (n \geqslant 2)$, 则 $I \subset \bigcup_n I_n$, 故
\[
m^* I \leqslant \sum_{n=1}^\infty l(I_n) = l(I_1) = b - a + 2\varepsilon.
\]
由 $\varepsilon$ 的任意性, $m^* I \leqslant b - a$.

\pause

\textbf{再证 $m^* I \geqslant b - a$}: 设 $\{I_n\}$ 是覆盖 $[a,b]$ 的任意开区间列. 由\textbf{有限覆盖定理}, 存在有限子覆盖 $I_{n_1}, \ldots, I_{n_k}$ 覆盖 $[a,b]$. 由区间长度的有限可加性,
\[
\sum_{n=1}^\infty l(I_n) \geqslant \sum_{j=1}^k l(I_{n_j}) \geqslant b - a.
\]
对上式取下确界得 $m^* I \geqslant b - a$.
\end{frame}

%------------------------------------------------

\begin{frame}{区间的外测度等于其长度}

综合两方面, $m^*[a,b] = b - a$. 对开区间、半开半闭区间的情形类似可证.

\pause

\begin{cor}
外测度是区间长度的真正推广: 对任意区间 $I$, $m^* I = l(I).$
\end{cor}

\pause

\begin{eg}
\begin{itemize}
\item $m^*[0,1] = 1$ (满足规范性).
\item $m^*(a, b) = b - a = m^*[a,b] = m^*[a,b) = m^*(a,b].$
\end{itemize}
\end{eg}

\pause

\begin{remark}
证明的关键在于``$\geqslant$''方向用到了\textbf{有限覆盖定理}, 这是 $\mathbb{R}$ 中闭区间的重要性质. 高维情形类似, 用闭矩体代替闭区间即可.
\end{remark}

\end{frame}

%------------------------------------------------

\section{内测度}

%------------------------------------------------

\begin{frame}{内测度的定义}

\begin{Def}[内测度 (inner measure)]
设 $E \subset [a,b]$ 为有界集. 定义 $E$ 的\alert{内测度 (inner measure)} 为
\[
m_* E = (b - a) - m^*([a,b] \setminus E).
\]
\end{Def}

\pause

\begin{itemize}
\item 几何直观: 用区间 $[a,b]$ 的长度减去``补集的外测度'', 从外部``夹''出 $E$ 的``内部大小''.
\item 内测度不依赖于包含区间的选取 (见下).
\end{itemize}

\pause

\begin{prop}[良定性]
若 $E \subset [a,b] \subset [c,d]$, 则
\[
(b-a) - m^*([a,b] \setminus E) = (d-c) - m^*([c,d] \setminus E).
\]
即 $m_* E$ 的定义与包含区间 $[a,b]$ 的选取无关.
\end{prop}

\pause

\startproof[bg=yellow, fg=red]
注意到 $[c,d] \setminus E = ([c,d] \setminus [a,b]) \cup ([a,b] \setminus E)$ 为不交并. 由次可加性及区间外测度 = 长度, 有
\[
m^*([c,d] \setminus E) \leqslant m^*([c,d] \setminus [a,b]) + m^*([a,b] \setminus E) = (d-c) - (b-a) + m^*([a,b] \setminus E).
\]
整理即得. 反向不等式同理.
\end{frame}

%------------------------------------------------

\begin{frame}{内测度的基本事实}

\begin{prop}
设 $E$ 为有界集, 则 $m_* E \leqslant m^* E.$
\end{prop}

\pause

\startproof[bg=yellow, fg=red]
不妨设 $E \subset [a,b]$. 对任意开区间列 $\{I_n\}$ 覆盖 $E$ 和开区间列 $\{J_k\}$ 覆盖 $[a,b] \setminus E$, 合并起来覆盖 $[a,b]$, 故
\[
\sum_n l(I_n) + \sum_k l(J_k) \geqslant b - a.
\]

\pause

对固定的 $\{I_n\}$, 在上式对 $\{J_k\}$ 取下确界, 得
\[
\sum_n l(I_n) + m^*([a,b] \setminus E) \geqslant b - a.
\]

\pause

再对 $\{I_n\}$ 取下确界, 即得
\[
m^* E + m^*([a,b] \setminus E) \geqslant b - a,
\]
亦即 $m^* E \geqslant (b-a) - m^*([a,b] \setminus E) = m_* E.$
\end{frame}

%------------------------------------------------

\begin{frame}{内测度的等价描述}

\begin{prop}[内测度的等价定义]
设 $E \subset [a,b]$ 为有界集, 则
\[
m_* E = \sup\{m F : F \subset E,\ F \text{ 为有界闭集}\}.
\]
\end{prop}

\pause

\startproof[bg=yellow, fg=red]
对任意有界闭集 $F \subset E$, 有 $[a,b] \setminus E \subset [a,b] \setminus F$, 由单调性
\[
m^*([a,b] \setminus E) \leqslant m^*([a,b] \setminus F).
\]
而 $[a,b] \setminus F$ 是开集, 利用开集的结构定理可证 $m^*([a,b] \setminus F) = (b-a) - m F$ (开集可测的结论, 详见 §3). 于是
\[
m_* E = (b-a) - m^*([a,b] \setminus E) \geqslant (b-a) - m^*([a,b] \setminus F) = m F.
\]
故 $m_* E \geqslant \sup\{m F : F \subset E,\ F \text{ 闭集}\}$. 反向不等式同理可证.
\end{frame}

%------------------------------------------------

\section{可测集的定义}

%------------------------------------------------

\begin{frame}{可测集的定义}

\begin{Def}[Lebesgue 可测集]
设 $E \subset \mathbb{R}$ 为有界集. 若 $m^* E = m_* E$, 则称 $E$ 为\alert{(Lebesgue) 可测集}. 此时称公共值为 $E$ 的\alert{Lebesgue 测度}, 记为 $m E$.
\end{Def}

\pause

所有可测集构成的集合称为\alert{可测集类}, 记为 $\mathscr{M}$.

\pause

\begin{remark}
\begin{itemize}
\item 对无界集 $E$, 定义 $E$ 可测 $\iff E \cap [-n, n]$ 可测 $\forall n \in \mathbb{N}$.
\item 本节只讨论有界集的情形.
\end{itemize}
\end{remark}

\pause

\begin{question}
哪些集合是可测的? 可测集类 $\mathscr{M}$ 有多大?
\end{question}

\end{frame}

%------------------------------------------------

\begin{frame}{可测集的初步例子}

\begin{prop}[区间可测]
任意区间 $I$ (开、闭、半开半闭) 都是可测的, 且 $m I = l(I).$
\end{prop}

\pause

\startproof[bg=yellow, fg=red]
设 $I = [a,b]$. 已证 $m^* I = b - a$. 又 $[a,b] \setminus I = \{a, b\}$ 是两点集, 外测度为 $0$, 故
\[
m_* I = (b-a) - m^*([a,b] \setminus I) = (b-a) - 0 = b - a = m^* I.
\]
故 $I$ 可测, $m I = b - a$.
\end{frame}

%------------------------------------------------

\begin{frame}{可测集的初步例子}

\begin{prop}[零测集可测]
若 $m^* E = 0$, 则 $E$ 可测, 且 $m E = 0$.
\end{prop}

\pause

\startproof[bg=yellow, fg=red]
不妨设 $E \subset [a,b]$. 由 $m^* E = 0$ 及单调性, $m^*([a,b] \setminus E) \leqslant b - a$, 故
\[
m_* E = (b-a) - m^*([a,b] \setminus E) \geqslant (b-a) - (b-a) = 0.
\]
另一方面, 由内测度的基本事实 $0 \leqslant m_* E \leqslant m^* E = 0$, 故 $m_* E = 0 = m^* E$, 即 $E$ 可测.
\end{frame}

%------------------------------------------------

\begin{frame}{可测集的等价刻画 (1)}

\begin{thm}[用开集 / 闭集逼近]
设 $E$ 为有界集. 则以下等价:
\begin{enumerate}
\item $E$ 可测;
\item $\forall\, \varepsilon > 0,$ 存在开集 $G \supset E$ 使得 $m^*(G \setminus E) < \varepsilon;$
\item $\forall\, \varepsilon > 0,$ 存在闭集 $F \subset E$ 使得 $m^*(E \setminus F) < \varepsilon.$
\end{enumerate}
\end{thm}

\pause

\startproof[bg=yellow, fg=red]
$(1) \Rightarrow (2)$: 由 $E$ 可测及外测度定义, 对 $\varepsilon > 0$, 存在开区间列 $\{I_n\}$ 覆盖 $E$ 使得 $\sum_n l(I_n) < m^* E + \varepsilon$. 令 $G = \bigcup_n I_n$, 则 $G$ 是开集, $G \supset E$, 且
\[
m^* E \leqslant m^* G \leqslant \sum_n l(I_n) < m^* E + \varepsilon.
\]
又 $G = E \cup (G \setminus E)$ 为不交并, 由 $E$ 可测及次可加性取等条件 (有限可加性, 见后), 有 $m G = m^* E + m^*(G \setminus E)$. 故 $m^*(G \setminus E) = m G - m^* E < \varepsilon$.
\end{frame}

%------------------------------------------------

\begin{frame}{可测集的等价刻画 (2)}

\startproof[bg=yellow, fg=red]
$(2) \Rightarrow (1)$: 对 $\varepsilon > 0$, 取开集 $G \supset E$ 使 $m^*(G \setminus E) < \varepsilon$. 由 $E \subset G$, 有 $m^* E \leqslant m^* G$. 又 $G = E \cup (G \setminus E)$, 由次可加性
\[
m^* G \leqslant m^* E + m^*(G \setminus E) < m^* E + \varepsilon.
\]
于是 $m^* E \leqslant m^* G < m^* E + \varepsilon$.

\pause

另一方面, 由内测度定义及 $G$ 的结构 (开集可测的结论), 有
\[
m_* E \geqslant m^* G - m^*(G \setminus E) > m^* G - \varepsilon > m^* E - 2\varepsilon.
\]

\pause
故 $m^* E - m_* E < 2\varepsilon$. 由 $\varepsilon$ 的任意性, $m^* E = m_* E$, 即 $E$ 可测.

\pause

$(1) \Leftrightarrow (3)$: 利用补集关系 $m_* E = (b-a) - m^*([a,b] \setminus E)$, 由 $(1) \Leftrightarrow (2)$ 对补集应用即得.
\end{frame}

%------------------------------------------------

\section{可测集关于有限运算的封闭性}

%------------------------------------------------

\begin{frame}{有限可加性}

\begin{prop}[有限可加性]
设 $E_1, E_2$ 是可测集且 $E_1 \cap E_2 = \emptyset$, 则 $E_1 \cup E_2$ 可测, 且
\[
m(E_1 \cup E_2) = m E_1 + m E_2.
\]
\end{prop}

\pause

\startproof[bg=yellow, fg=red]
由 $E_1$ 可测及等价刻画, 对 $\varepsilon > 0$, 存在开集 $G_1 \supset E_1$ 使 $m^*(G_1 \setminus E_1) < \varepsilon$. 同理存在开集 $G_2 \supset E_2$ 使 $m^*(G_2 \setminus E_2) < \varepsilon$.

\pause

(详细证明需利用 $E_1 \cap E_2 = \emptyset$ 分离 $G_1, G_2$, 再用外测度次可加性和内测度定义验证 $m^*(E_1 \cup E_2) = m_*(E_1 \cup E_2) = m E_1 + m E_2$.)
\end{frame}

%------------------------------------------------

\begin{frame}{差运算}

\begin{prop}[差运算]
设 $E_1, E_2$ 可测且 $E_1 \subset E_2$, 则 $E_2 \setminus E_1$ 可测, 且
\[
m(E_2 \setminus E_1) = m E_2 - m E_1.
\]
\end{prop}

\pause

\startproof[bg=yellow, fg=red]
由 $E_2 = E_1 \cup (E_2 \setminus E_1)$ 为不交并, 利用有限可加性即得.
\end{frame}

%------------------------------------------------

\begin{frame}{小结}

\begin{itemize}
\item \textbf{外测度} $m^*$: 对所有有界集定义, 满足非负性、单调性、次可加性、平移不变性.
\item \textbf{内测度} $m_*$: 对有界集定义, 满足 $m_* E \leqslant m^* E$.
\item \textbf{可测集}: $m^* E = m_* E$. 区间、零测集可测.
\item \textbf{可测集关于有限并、差运算封闭}.
\end{itemize}

\pause

\begin{question}
\begin{itemize}
\item 可测集关于\textbf{可列并}是否封闭? (是, 但证明不简单, 见 §3.)
\item 开集是否可测? Borel 集呢? (是, 见 §3.)
\item 是否存在\textbf{不可测集}? (是, 见补充材料.)
\end{itemize}
\end{question}

\end{frame}

%------------------------------------------------

\end{document}
